% Chapter 5

\chapter{General Discussion, Limitations, and Future Directions} % Main chapter title

\label{Chapter5} % For referencing the chapter elsewhere, use \ref{Chapter1} 

%----------------------------------------------------------------------------------------

% Define some commands to keep the formatting separated from the content 
% \newcommand{\keyword}[1]{\textbf{#1}}
% \newcommand{\tabhead}[1]{\textbf{#1}}
% \newcommand{\code}[1]{\texttt{#1}}
% \newcommand{\file}[1]{\texttt{\bfseries#1}}
% \newcommand{\option}[1]{\texttt{\itshape#1}}

%----------------------------------------------------------------------------------------

\section{General Overview of the Thesis}

The central objective of this thesis is to address the long-standing challenge of inferring GRNs from time series gene expression data. This problem is fundamentally difficult due to the high dimensionality of biological systems, the limited number of observed time points, the presence of noise, and the intrinsic multimodality of network dynamics. Rather than focusing solely on improving inference algorithms, this thesis adopts a broader perspective, emphasizing the critical roles of data representation and data availability in enhancing GRN inference performance.

Chapter ~\ref{Chapter2}  introduces a multiple-level discretization framework for network inference. Unlike most existing discretization-based methods that rely on a fixed Boolean representation using two discrete states, this chapter proposes a flexible discretization strategy in which each gene can be represented using either two or three discrete levels, depending on the distribution of its expression values. By avoiding a rigid discretization scheme and allowing gene-specific discretization levels, the proposed approach preserves more biologically meaningful information while maintaining computational efficiency. This adaptive representation significantly improves both inference accuracy and runtime performance, demonstrating that discretization, when properly designed, can be a powerful tool rather than a limiting factor in GRN inference.

Building upon the importance of data representation, Chapter ~\ref{Chapter3} proposes a novel representation learning framework that jointly encodes gene expression dynamics and network structural information. Instead of directly inferring regulatory relationships from raw or discretized expression profiles, this framework transforms the GRN inference problem into a supervised learning task. By constructing informative feature representations that capture both temporal expression patterns and structural dependencies, standard machine learning classifiers can be effectively applied for network prediction. This chapter highlights that learning suitable representations can bridge the gap between biological network inference and modern machine learning techniques, enabling more scalable and flexible inference models.

While Chapters ~\ref{Chapter2} and ~\ref{Chapter3} focus on improving inference performance through better representations, Chapter ~\ref{Chapter4} addresses a complementary but equally important challenge: data scarcity. In many practical scenarios, the number of observed time points is severely limited, which substantially degrades inference accuracy. To mitigate this issue, Chapter ~\ref{Chapter4} proposes a data synthesis approach that generates additional time-series samples using an autoencoder-based model. By augmenting the original dataset with synthesized expression trajectories, the proposed method enhances the robustness and reliability of network inference under low-data conditions. This chapter demonstrates that improving data quantity, when done carefully, can significantly boost inference performance without requiring additional experimental measurements.

Taken together, the contributions of this thesis demonstrate that accurate GRN inference is not solely determined by the choice of inference algorithm. Instead, it critically depends on how gene expression data are represented, how structural information is incorporated, and how data limitations are addressed. By systematically exploring these aspects across Chapters ~\ref{Chapter2} to ~\ref{Chapter4}, this thesis provides a comprehensive framework for improving GRN inference from time-series gene expression data.


\section{Cross-Method Discussion and Integration of Findings}

\subsection{Data Representation as a Core Factor in GRN Inference}

Data representation plays a central and often underappreciated role in GRN inference. While a large body of existing studies has focused on designing increasingly sophisticated inference algorithms, comparatively less attention has been paid to how gene expression data are encoded prior to inference. In practice, the quality of the inferred network is strongly influenced by the representation of the input data, as it directly determines the amount of biological information preserved, the robustness to noise, and the computational tractability of the inference process.

Traditional GRN inference methods can be broadly categorized according to the form of data representation they employ. Continuous valued representations retain the full resolution of gene expression measurements and are theoretically capable of capturing subtle regulatory effects. However, they are highly sensitive to noise, require large numbers of samples, and often lead to prohibitively high computational complexity, especially for large-scale networks. On the other hand, Boolean representations simplify gene expression into binary states, enabling fast inference and facilitating the modeling of gene dynamics. Despite these advantages, Boolean discretization inevitably leads to information loss and may oversimplify complex regulatory behaviors, particularly when genes exhibit intermediate or graded expression levels.

The findings of this thesis demonstrate that neither extreme fully continuous nor strictly Boolean representations is optimal for GRN inference. Instead, an intermediate representation that balances information preservation and computational efficiency is more suitable. In Chapter ~\ref{Chapter2}, we introduced a multiple level discretization strategy in which genes are discretized into either two or three levels depending on the distribution of their expression values. This adaptive discretization avoids enforcing a uniform representation across all genes and acknowledges the inherent heterogeneity of gene expression dynamics. As a result, biologically meaningful variations are preserved for genes that require finer resolution, while computational efficiency is maintained for genes that can be adequately represented using fewer states.

Furthermore, Chapter ~\ref{Chapter3} extends the notion of data representation beyond discretization by proposing a representation learning framework that integrates gene expression values with structural information. This approach reframes GRN inference as a classification problem, where the success of inference is heavily dependent on how effectively the underlying biological relationships are encoded into feature representations. The improved performance observed in this framework highlights that appropriate representations can enable the use of powerful machine learning models, even when the original biological data are noisy or limited.

Chapter ~\ref{Chapter4} further reinforces the importance of representation by addressing data sparsity through synthetic data generation. By augmenting time-series data using an autoencoder based model, the effective representation of gene dynamics is enriched in the temporal dimension. This not only improves inference accuracy but also stabilizes the learning process when the number of observed time points is small. Importantly, the success of data augmentation depends on the quality of the learned latent representation, again underscoring representation as a key factor in GRN inference.

Taken together, the results presented in this thesis suggest that data representation should be considered a foundational component of GRN inference, rather than a mere preprocessing step. Adaptive discretization, structured feature encoding, and learned latent representations collectively contribute to more accurate, scalable, and biologically meaningful network inference. Future GRN inference studies are therefore likely to benefit from placing greater emphasis on representation design alongside algorithmic development.

\subsection{Structural Accuracy vs Dynamic Accuracy}
Structural accuracy and dynamic accuracy capture two fundamentally different aspects of gene regulatory network inference. Structural accuracy evaluates how well the inferred network topology matches the underlying goldstandard structure, whereas dynamic accuracy assesses the ability of the inferred network to reproduce observed gene expression trajectories. In complex biological systems, these two objectives are not necessarily aligned. Distinct network structures can generate identical or highly similar dynamic behaviors, leading to a multimodal inference landscape.

Our experimental results consistently demonstrate that dynamic accuracy remains high even when structural accuracy degrades, particularly as network size and the number of regulatory inputs increase. This observation highlights the intrinsic ambiguity of GRN inference from timeseries data and suggests that functional validity may be preserved despite structural uncertainty. Consequently, evaluating network inference methods solely based on structural metrics may be insufficient. Incorporating dynamic accuracy provides a more comprehensive assessment of predictive performance and biological relevance, especially for largescale and densely connected networks.
\subsection{Scalability and Network Complexity}
Scalability remains a critical challenge in gene regulatory network inference, particularly as network size and regulatory complexity increase. As the number of genes grows, the combinatorial search space for potential regulatory interactions expands rapidly, resulting in increased computational cost and inference difficulty. Moreover, higher network density, reflected by a larger number of incoming regulatory links, further exacerbates this challenge by introducing additional ambiguity in identifying true regulatory relationships.

Our results indicate that MIDNI maintains robust dynamic inference performance even in largescale and densely connected networks, despite a gradual decline in structural accuracy. This behavior suggests that the proposed multi-level discretization and information theoretic feature selection strategy effectively mitigates the impact of network complexity. While increased computational cost is unavoidable in large networks, the observed trade-off between scalability and inference accuracy remains favorable, supporting the applicability of MIDNI to realistic biological systems where both network size and regulatory complexity are substantial.

\section{Limitations of the Proposed Approaches}
\subsection{Multimodality of the Network Inference Problem}
A fundamental limitation of gene regulatory network inference lies in the multimodality of the solution space, where multiple distinct network structures are capable of producing similar or even identical dynamic behaviors. This characteristic becomes particularly prominent in large and densely connected networks, in which the number of plausible regulatory configurations increases combinatorially. As a result, optimizing for dynamic accuracy alone does not guarantee the recovery of the true underlying network structure. In this study, MIDNI frequently identified networks that closely reproduced the observed gene expression dynamics while differing from the gold standard topology. This phenomenon reflects an inherent ambiguity in time series based inference rather than a deficiency of the proposed method. Consequently, structural inference accuracy is intrinsically constrained by the limited observability of gene expression data. Addressing this limitation likely requires the integration of additional biological constraints, prior knowledge, or experimental perturbation data to reduce the space of equivalent solutions.
\subsection{Limitations of Mutual Information-Based Selection}
Another important limitation of the proposed framework arises from the use of mutual information as the primary criterion for selecting candidate regulatory genes. Mutual information is effective in capturing statistical dependencies between gene expression profiles, including nonlinear relationships, but it does not explicitly distinguish between direct and indirect interactions. As a consequence, genes that influence a target gene through intermediate regulators may exhibit high mutual information values and be incorrectly selected as direct regulators. This issue becomes more pronounced in dense networks, where indirect dependencies are common and difficult to disentangle using pairwise or approximate multivariate measures alone. In addition, mutual information does not encode causal directionality, which further complicates accurate regulator identification. While the SWAP subroutine partially alleviates this problem by refining candidate sets based on dynamic consistency, the reliance on mutual information inherently limits structural precision. Incorporating causal inference mechanisms or biological prior knowledge may help mitigate this limitation in future work.
\subsection{Limitations in Direction and Sign Determination}
A further limitation of the proposed approaches concerns the determination of regulatory direction and interaction sign within inferred gene regulatory networks. While the models effectively identify the presence of regulatory relationships, accurately inferring whether a regulation is activating or inhibiting remains challenging, particularly under noisy or sparsely sampled time series conditions. Directionality is often inferred indirectly from temporal ordering or model predictions, yet such cues can be ambiguous when gene expression dynamics exhibit delays, feedback loops, or synchronized behavior. Moreover, classification based frameworks tend to prioritize structural prediction accuracy, which may lead to correct edge detection but incorrect assignment of regulatory sign. This limitation affects the biological interpretability of the inferred networks, since activation and repression convey fundamentally different functional meanings. Although increasing temporal resolution and integrating dynamic modeling can partially improve sign inference, the inherent ambiguity in expression based data constrains reliable direction and sign determination without additional biological constraints or intervention data.
\subsection{Limitations of Synthetic Data Generation}
Despite the effectiveness of synthetic data generation in alleviating data scarcity, this strategy introduces several inherent limitations. The quality of generated samples strongly depends on the amount and diversity of the available original data, and when the initial training set is limited, the learned generative model may fail to capture meaningful temporal patterns, leading to degraded inference performance. In addition, the number of synthesized time points must be carefully controlled, since the generated data at one stage are often used as input for subsequent generation steps, which results in the accumulation of modeling errors over time. As the generation process continues, these errors can propagate and amplify, causing the synthetic data to gradually deviate from realistic gene expression dynamics. Furthermore, synthetic samples may introduce bias by over reinforcing dominant patterns while suppressing rare but biologically important behaviors. Consequently, excessive reliance on synthetic data can compromise both the accuracy and biological validity of the inferred networks.
\section{Future Research Directions}
\subsection{Beyond Greedy Search: Advanced Optimization Strategies}
The inference framework proposed in this thesis relies on greedy search strategies in the MIFS and SWAP subroutines, which provide computational efficiency and stable performance but are inherently limited when exploring complex and high dimensional solution spaces. Greedy algorithms tend to converge to locally optimal solutions, particularly in multimodal inference problems where multiple regulatory configurations can produce similar dynamic behaviors. To overcome this limitation, future research should investigate advanced optimization strategies that allow broader exploration of the search space while preserving computational feasibility. Potential directions include metaheuristic approaches such as genetic algorithms, simulated annealing, and particle swarm optimization, which have demonstrated effectiveness in navigating multimodal landscapes. In addition, hybrid optimization frameworks that combine greedy initialization with global search refinement may offer a favorable balance between efficiency and accuracy. Incorporating such strategies is expected to improve both structural inference and robustness in large scale gene regulatory networks.
\subsection{Causal and Mechanistic Inference}
Most data driven gene regulatory network inference methods, including the approaches proposed in this thesis, primarily focus on statistical dependencies rather than true causal and mechanistic relationships among genes. Although inferred interactions may exhibit high structural or dynamical accuracy, they do not necessarily reflect direct regulatory effects or underlying biological mechanisms. This limitation becomes more pronounced when genes are indirectly connected through hidden regulators or unobserved confounding factors. Future research should aim to integrate causal inference frameworks that explicitly model intervention effects and temporal precedence, allowing a clearer distinction between correlation and causation. In addition, incorporating mechanistic constraints derived from transcription factor binding, signaling pathways, or kinetic models may enhance biological interpretability. A hybrid framework that combines machine learning based inference with causal reasoning and mechanistic knowledge could lead to more reliable and biologically meaningful gene regulatory networks.
\subsection{Integration of Prior Biological Knowledge}
While the proposed methods rely on data driven learning from gene expression profiles, incorporating prior biological knowledge represents an important direction for improving the reliability and interpretability of inferred gene regulatory networks. Purely data driven approaches may suffer from false positive interactions, especially when expression data are noisy, sparse, or limited in temporal resolution. Integrating information from biological databases such as known transcription factor target relationships, protein DNA binding data, or curated signaling pathways can help constrain the search space and guide the inference process toward biologically plausible solutions. Prior knowledge can be incorporated in various forms, including regularization terms, probabilistic priors, or post inference filtering strategies. Such integration not only has the potential to improve structural accuracy but also enhances biological credibility, enabling inferred networks to better reflect true regulatory mechanisms rather than statistical artifacts.
\subsection{Extension to Other Data Types}
All three methods proposed in this thesis are developed and evaluated based on time series gene expression data, which provide valuable temporal information for capturing regulatory dynamics. However, another widely available and biologically relevant data type is steady state gene expression data, which reflect equilibrium conditions under different perturbations or environmental settings. Steady state data are often easier to obtain and can complement time series observations, especially when the number of observed time points is limited. A promising future direction is to design an integrated data representation framework that jointly exploits both time series and steady state expression data. By combining dynamic temporal patterns with steady state dependencies, such a unified representation has the potential to enhance network inference robustness and accuracy, particularly in scenarios where either data type alone is insufficient.

\section{Concluding Remarks}
This thesis investigated the problem of gene regulatory network inference from gene expression data, with a particular focus on improving inference performance through effective data representation and learning strategies. Across the three proposed frameworks, we demonstrated that how gene expression data are represented plays a critical role in determining both structural and dynamical inference accuracy. By moving beyond conventional Boolean representations and raw real valued data, this work introduced flexible multi level discretization, representation learning that integrates expression and structural information, and synthetic data generation to address data scarcity in time series settings.

Extensive experiments on benchmark datasets showed that the proposed methods consistently outperformed existing approaches, especially for large scale and complex networks. These results highlight the importance of adapting inference strategies to data characteristics rather than relying on fixed modeling assumptions. Although several limitations remain, including optimization challenges and dependency on statistical measures, this thesis provides a unified perspective on network inference that bridges data representation, machine learning, and biological relevance.

Overall, this work contributes methodological insights and practical frameworks that can serve as a foundation for future advances in robust and scalable gene regulatory network inference.





